\begin{tikzpicture}
\begin{semilogxaxis}[
    title={Diagrama de Bode de Magnitud},
    xlabel={Frecuencia $\omega$ (rad/s)},
    ylabel={Magnitud (dB)},
    grid=both,
    grid style={line width=.1pt, draw=gray!30},
    major grid style={line width=.2pt, draw=gray!60},
    xmin=1e0, xmax=1e5,
    ymin=-40, ymax=40,
    width=0.9\columnwidth,
    height=8cm,
    legend style={
    nodes={scale=0.7, transform shape}
}
]

\addplot[thick, blue] coordinates {
    (1, -6.02)
    (40, 26.02)
    (160, 26.02)
    (1e5, -29.90)
};
\addlegendentry{Teórico (Asintótico)}

\addplot[mark=*, mark size=1.5pt, red, forget plot] coordinates {(40, 26.02)};
\node[anchor=south east, font=\small, text=black] at (axis cs:40, 26.02) {$\omega_{p1} = 40$};

\addplot[mark=*, mark size=1.5pt, red, forget plot] coordinates {(160, 26.02)};
\node[anchor=south west, font=\small, text=black] at (axis cs:160, 26.02) {$\omega_{p2} = 160$};

% Datos de LTSpice
\addplot[thick, dashed, orange] table [
    skip first n=1,
    x expr=\thisrowno{0} * 2 * pi,
    y index=1
] {sources/circuito3_cleaned.txt};
\addlegendentry{Simulación LTSpice}

\end{semilogxaxis}
\end{tikzpicture}
